\documentclass[aspectratio=169,11pt]{beamer}
\usepackage{../phanes-beamer}
\renewcommand{\ModuleID}{07}
\title{RAG for regulatory review}
\subtitle{Retrieve traceable public evidence and bound the interpretation}
\author{PHANES Initiative}
\institute{The University of Texas at Austin}
\date{September 12, 2026}
\begin{document}
\begin{frame}[plain]
\titlepage
\end{frame}
\begin{frame}[t,fragile]{Retrieval supports the reviewer's judgment}
\fontsize{11.5}{14}\selectfont
\begin{itemize}\setlength{\itemsep}{2mm}
\item Find relevant public requirements and guidance.
\item Map a supplied statement to exact source passages.
\item Distinguish retrieved evidence from interpretation.
\item Flag missing evidence and uncertain applicability for review.
\end{itemize}
\end{frame}
\begin{frame}[t,fragile]{The regulatory review pipeline}
\vspace{-1mm}
\begin{center}
\begin{tikzpicture}[
  node distance=4mm and 8mm,
  every node/.style={draw=UTorange,rounded corners,align=center,font=\scriptsize,minimum height=6mm,minimum width=38mm},
  >={Stealth}
]
\node (corpus) {Public corpus:\\10 CFR, Reg Guides, SRP};
\node[below=of corpus] (manifest) {Corpus manifest:\\URL, date, SHA-256};
\node[below=of manifest] (chunk) {Clause-aware\\chunking};

\node[right=of corpus] (index) {Hybrid index:\\lexical + embedding};
\node[below=16mm of index] (app) {Application statement\\to evaluate};

\node[right=of index] (retrieval) {Candidate retrieval\\\& reranking};
\node[below=of retrieval] (claim) {Claim-to-evidence\\record};
\node[below=of claim] (reviewer) {Human reviewer:\\judgment \& gaps};

\draw[->] (corpus) -- (manifest);
\draw[->] (manifest) -- (chunk);
\draw[->] (chunk.east) to[bend right=15] (index.south west);
\draw[->] (index) -- (retrieval);
\draw[->] (app.north) -- (retrieval.south west);
\draw[->] (retrieval) -- (claim);
\draw[->] (claim) -- (reviewer);
\end{tikzpicture}
\end{center}
\vspace{-1mm}
\fontsize{10.5}{12.5}\selectfont
\begin{itemize}\setlength{\itemsep}{1.5mm}
\item Ingestion is offline and provenance-tracked: every chunk links to its manifest hash and exact text span.
\item Retrieval produces candidate evidence; the human reviewer retains authority over regulatory findings.
\end{itemize}
\end{frame}
\begin{frame}[t,fragile]{Preserve the kind of authority}
\fontsize{10.5}{12.5}\selectfont\setlength{\tabcolsep}{4pt}\renewcommand{\arraystretch}{1.15}
\begin{tabular}{>{\raggedright\arraybackslash}p{\dimexpr 0.27\linewidth-4.32pt\relax}>{\raggedright\arraybackslash}p{\dimexpr 0.73\linewidth-11.68pt\relax}}
\toprule
\textbf{Document type} & \textbf{How to treat it} \\ \midrule
Regulation & Read the applicable rule and scope \\[2pt]
Regulatory guide & Identify the described acceptable approach \\[2pt]
Standard review plan & Understand staff review guidance \\[2pt]
Application statement & Treat as a claim needing evidence \\[2pt]
Technical report & Assess relevance, method and limitations \\[2pt]
\bottomrule\end{tabular}\par\vspace{2mm}
\vfill{\fontsize{6.5}{8}\selectfont\color{UTgray} Sources: rg, srp. See references and instructor notes.\par}
\end{frame}
\begin{frame}[t,fragile]{Create a corpus manifest before indexing}
\fontsize{11.5}{14}\selectfont
\begin{itemize}\setlength{\itemsep}{2mm}
\item Record source URL, title, document identifier, and revision/date.
\item Retain the downloaded artifact and its hash.
\item Record section and PDF page separately when they differ.
\item Keep public status and extraction limitations visible.
\end{itemize}
\end{frame}
\begin{frame}[t,fragile]{Chunk by meaning and citation boundaries}
\fontsize{11.5}{14}\selectfont
\begin{itemize}\setlength{\itemsep}{2mm}
\item Keep a requirement with its qualifiers and definitions.
\item Preserve section headings and document metadata.
\item Handle tables, footnotes, and OCR errors explicitly.
\item Avoid splitting "shall" statements from their conditions.
\end{itemize}
\end{frame}
\begin{frame}[t,fragile]{Compare retrieval methods on the same queries}
\fontsize{10.5}{12.5}\selectfont\setlength{\tabcolsep}{4pt}\renewcommand{\arraystretch}{1.15}
\begin{tabular}{>{\raggedright\arraybackslash}p{\dimexpr 0.25\linewidth-6.0pt\relax}>{\raggedright\arraybackslash}p{\dimexpr 0.35\linewidth-8.4pt\relax}>{\raggedright\arraybackslash}p{\dimexpr 0.4\linewidth-9.6pt\relax}}
\toprule
\textbf{Method} & \textbf{Useful behavior} & \textbf{Failure to watch} \\ \midrule
Lexical search & Exact terms and identifiers & Missed paraphrases \\[2pt]
Embedding search & Semantic similarity & Similar but inapplicable text \\[2pt]
Hybrid retrieval & Combines candidate signals & Poor ranking/calibration \\[2pt]
Metadata filtering & Restricts revision or scope & Incorrect metadata exclusion \\[2pt]
\bottomrule\end{tabular}\par\vspace{2mm}
\end{frame}
\begin{frame}[t,fragile]{A high-scoring passage is the wrong revision}
\fontsize{11.5}{14}\selectfont
\begin{itemize}\setlength{\itemsep}{2mm}
\item Two passages use nearly identical language.
\item The top result belongs to a superseded document.
\item Which metadata should constrain retrieval?
\item How should the report expose the selected revision?
\end{itemize}
\end{frame}
\begin{frame}[t,fragile]{One worked approach: the superseded passage}
\fontsize{11.5}{14}\selectfont
\begin{itemize}\setlength{\itemsep}{2mm}
\item Constraining metadata: document identity, revision or date, and applicability to the reviewed case; the query is filtered so both revisions can appear and be compared.
\item Selection rule: the current applicable revision is used for the answer, and the superseded passage is retained in the record, not silently dropped.
\item Report exposure: the citation names the selected revision and date, and a note marks the other passage as superseded with the same document identity.
\item Caution: the latest revision is not always the correct one for the reviewed application; applicability decides, and the choice is stated in the claim-to-evidence record.
\end{itemize}
\end{frame}
\begin{frame}[t,fragile]{Build a claim-to-evidence record}
\fontsize{11.5}{14}\selectfont
\begin{itemize}\setlength{\itemsep}{2mm}
\item Application statement and the question under review.
\item Source identifier, revision, section, and exact supporting span.
\item Interpretation with assumptions stated separately.
\item Status: supported, contradicted, or insufficient evidence.
\end{itemize}
\end{frame}
\begin{frame}[t,fragile]{Make an unresolved finding explicit}
\begin{lstlisting}
{
  "statement_id": "DEMO-03",
  "status": "insufficient_evidence",
  "evidence_ids": [],
  "interpretation": "The searched corpus does not resolve this claim.",
  "next_action": "Human reviewer identifies the missing source."
}
\end{lstlisting}
\end{frame}
\begin{frame}[t,fragile]{Keep an evidence span and interpretation separate}
\begin{lstlisting}
{"claim_id":"DEMO-12", "corpus_version":"public-v1",
 "source_id":"DOC-04", "revision":"retained-edition",
 "section":"reviewed-section", "page":12,
 "span_start":1042, "span_end":1189,
 "evidence_text":"[verbatim text retained in corpus]",
 "interpretation":"[reviewer-visible inference]",
 "status":"needs_review"}
\end{lstlisting}
\fontsize{10.5}{12.5}\selectfont
\begin{itemize}\setlength{\itemsep}{2mm}
\item Schema example only: these placeholders are not a real regulatory citation.
\item Validate locators against the retained source before accepting an answer.
\end{itemize}
\end{frame}
\begin{frame}[t,fragile]{Check citations independently of fluency}
\fontsize{11.5}{14}\selectfont
\begin{itemize}\setlength{\itemsep}{2mm}
\item Does the cited document exist in the retained corpus?
\item Does the locator reach the quoted or paraphrased passage?
\item Does the passage support the specific claim?
\item Are qualifiers, revision, and applicability preserved?
\end{itemize}
\end{frame}
\begin{frame}[t,fragile]{Calculate retrieval precision and recall}
\fontsize{11.5}{14}\selectfont
\begin{itemize}\setlength{\itemsep}{2mm}
\item A labeled query has four relevant passages.
\item The system retrieves five passages; three are relevant.
\item Calculate precision and recall for this retrieval set.
\item Which missing passage would matter most to the reviewer?
\end{itemize}
\end{frame}
\begin{frame}[t,fragile]{Worked solution: precision and recall}
\fontsize{11.5}{14}\selectfont
\begin{itemize}\setlength{\itemsep}{2mm}
\item Precision: $3/5 = 0.60$; three of the five retrieved passages are relevant.
\item Recall: $3/4 = 0.75$; three of the four relevant passages were retrieved.
\item The missing passage is the one whose absence changes the review outcome, for example a key exception or a scope qualifier, not merely the highest-ranked absent item.
\item The counts are invented for the exercise; small labeled sets need reviewer agreement before the metrics support a decision.
\end{itemize}
\end{frame}
\begin{frame}[t,fragile]{Compare two retrieval configurations}
\fontsize{11.5}{14}\selectfont
\begin{itemize}\setlength{\itemsep}{2mm}
\item Illustrative labeled set: 10 queries with 40 relevant spans in total.
\item System A retrieves 50 spans, 30 relevant: micro precision 0.60 and micro recall 0.75.
\item System B retrieves 80 spans, 36 relevant: micro precision 0.45 and micro recall 0.90.
\item Choose based on review cost and missed-requirement consequence; also inspect per-query failures and applicability.
\end{itemize}
\end{frame}
\begin{frame}[t,fragile]{Evaluate retrieval and answers separately}
\fontsize{10.5}{12.5}\selectfont\setlength{\tabcolsep}{4pt}\renewcommand{\arraystretch}{1.15}
\begin{tabular}{>{\raggedright\arraybackslash}p{\dimexpr 0.27\linewidth-4.32pt\relax}>{\raggedright\arraybackslash}p{\dimexpr 0.73\linewidth-11.68pt\relax}}
\toprule
\textbf{Stage} & \textbf{Metric or check} \\ \midrule
Retrieval & Recall@k, precision@k, relevant revision \\[2pt]
Citation & Valid locator and supporting span \\[2pt]
Answer & Unsupported claims and missed qualifications \\[2pt]
Review finding & Missed issue and false-positive rate \\[2pt]
Boundary & Only approved corpus and endpoint paths \\[2pt]
\bottomrule\end{tabular}\par\vspace{2mm}
\end{frame}
\begin{frame}[t,fragile]{Use adversarial but public test cases}
\fontsize{11.5}{14}\selectfont
\begin{itemize}\setlength{\itemsep}{2mm}
\item A query uses a synonym instead of the exact regulatory term.
\item A scanned table contains an extraction error.
\item A passage includes a relevant exception in the next paragraph.
\item A document contains instructions unrelated to the review task.
\end{itemize}
\end{frame}
\begin{frame}[t,fragile]{Lab: build a small review assistant}
\fontsize{11.5}{14}\selectfont
\begin{itemize}\setlength{\itemsep}{2mm}
\item Curate a public corpus and retain its manifest.
\item Implement a CPU lexical baseline; add approved embeddings optionally.
\item Use OpenCode to build cited statement-to-evidence mapping.
\item Evaluate retrieval, citation correctness, and unsupported findings.
\end{itemize}
\end{frame}
\begin{frame}[t,fragile]{Sample submission: the review assistant}
\fontsize{11.5}{14}\selectfont
\begin{itemize}\setlength{\itemsep}{2mm}
\item Corpus and manifest: the public documents used, the extraction method, the corpus version, and the instructor-chosen statement labels.
\item Baseline and evaluation: the CPU lexical retriever with labeled queries, and per-query precision/recall with the ambiguous cases documented.
\item Statement-to-evidence mapping: each claim carries a source identifier, revision, section, and the supporting span with a supported or insufficient status.
\item Findings: the unsupported claims are listed explicitly, and the citation checks separate valid locators from entailment.
\end{itemize}
\end{frame}
\begin{frame}[t,fragile]{Exit ticket: when should the assistant stop?}
\fontsize{11.5}{14}\selectfont
\begin{itemize}\setlength{\itemsep}{2mm}
\item No supporting passage in the searched corpus.
\item Conflicting or unclear document applicability.
\item Unreliable extraction or missing qualifiers.
\item A conclusion that requires a regulator or qualified reviewer.
\end{itemize}
\end{frame}
\begin{frame}[t,fragile]{A complete answer: when to stop}
\fontsize{11.5}{14}\selectfont
\begin{itemize}\setlength{\itemsep}{2mm}
\item No supporting passage: status is insufficient evidence, and the next action is to have a human reviewer identify the missing source.
\item Conflicting applicability: status is needs review, and the record names both revisions and the question that distinguishes them.
\item Unreliable extraction: status is unverified, and the claim is held until a trusted extraction or the original page confirms it.
\item Regulator-level conclusion: the assistant stops at a status record and a next action for the qualified reviewer; it does not render the compliance verdict.
\end{itemize}
\end{frame}
\begin{frame}[t]{References 1}
\fontsize{9}{12}\selectfont\begin{itemize}\setlength{\itemsep}{3mm}
\item \textbf{curriculum}: User-supplied PHANES curriculum: mod/01 through mod/08/README.md.
\item \textbf{colors}: \href{https://umac.utexas.edu/brand-center/colors/}{umac.utexas.edu/brand-center/colors}
\item \textbf{rg}: \href{https://www.nrc.gov/reading-rm/doc-collections/reg-guides/index}{nrc.gov/reading-rm/doc-collections/reg-guides/index}
\item \textbf{srp}: \href{https://www.nrc.gov/reading-rm/doc-collections/nuregs/staff/sr0800/index}{nrc.gov/reading-rm/doc-collections/nuregs/staff/sr0800/index}
\end{itemize}\end{frame}
\end{document}